\documentclass[12pt]{article}
\usepackage[T1]{fontenc}
\usepackage[utf8]{inputenc}
\usepackage{lmodern}
\usepackage{textcomp}
\usepackage{enumitem}
\usepackage{geometry}
\geometry{margin=1in}
\begin{document}
\begin{center}
Answer Key - Sociology - Undergraduate Level Assessment 10 \
\small (Answers are ordered exactly as in the question paper.)
\end{center}

\section*{Multiple Choice}
\begin{enumerate}[label=\arabic*.]
\item \textbf{Answer:} B
\\ \textit{Options:} A) the household; B) the office; C) the global village; D) the nation state
\\ \textit{Note:} The office is not recognized as a level of society because it represents a specific organizational context rather than a broader social structure, which typically includes levels such as individuals, groups, institutions, and societies.
\item \textbf{Answer:} A
\\ \textit{Options:} A) we internalise and take on social roles from a pre-existing framework; B) we create and negotiate our roles through interaction with others; C) social roles are not fixed or stable but fluid and pluralistic; D) roles have to be learned to suppress unconscious motivations
\\ \textit{Note:} Role-learning theory posits that individuals learn and internalize their social roles through interactions within a pre-existing societal framework, which shapes their behaviors and expectations in various contexts.
\item \textbf{Answer:} C
\\ \textit{Options:} A) government 'white paper'; B) confidential medical records; C) household account book; D) the shares register of a business
\\ \textit{Note:} A household account book is a personal document that records private financial transactions and is typically kept confidential, making it a closed-access document.
\item \textbf{Answer:} D
\\ \textit{Options:} A) Prime Minister, Archbishop of York, Viscounts of England; B) Marquesses of England, Earls of Great Britain, King's Brothers; C) Esquires, Serjeants of Law, Dukes' Eldest Sons; D) King's Grandsons, Lord High Treasurer, Companions of the Bath
\\ \textit{Note:} The correct order reflects the traditional British social hierarchy, where the King's Grandsons hold the highest status as royal family members, followed by the Lord High Treasurer as a high-ranking official, and finally, the Companions of the Bath, who are recognized for their service but rank lower in the hierarchy.
\item \textbf{Answer:} A
\\ \textit{Options:} A) sensitive, caring, and emotional; B) laddish, aggressive, and violent; C) a strong and dependable breadwinner; D) openly bisexual and proud of it
\\ \textit{Note:} The correct answer highlights the shifting societal expectations of masculinity during the 1980 s, where the 'new man' was characterized by emotional openness and a nurturing demeanor, contrasting with traditional notions of masculinity that emphasized stoicism and toughness.
\end{enumerate}

\section*{True / False}
\begin{enumerate}[label=\arabic*.]
\setcounter{enumi}{5}
\item \textbf{Answer:} True
\\ \textit{Options:} A) True; B) False
\\ \textit{Note:} The introduction of market principles in the 1980 s often resulted in increased state regulation as governments sought to ensure standards, accountability, and quality in public services through mechanisms like national testing and inspections to balance market efficiency with social equity.
\item \textbf{Answer:} True
\\ \textit{Options:} A) True; B) False
\\ \textit{Note:} In contemporary societies, social institutions such as family, education, religion, and economy are interconnected and operate through specialized roles and practices that collectively shape social behavior and structure.
\item \textbf{Answer:} True
\\ \textit{Options:} A) True; B) False
\\ \textit{Note:} The nuclear family is defined as a family unit consisting of two generations, specifically parents and their biological or adopted children, living together in a single household.
\item \textbf{Answer:} True
\\ \textit{Options:} A) True; B) False
\\ \textit{Note:} The 1944 Education Act, also known as the Butler Act, established a framework for free secondary education for all children in England and Wales, ensuring that education was accessible to every child regardless of their socio-economic background.
\item \textbf{Answer:} True
\\ \textit{Options:} A) True; B) False
\\ \textit{Note:} The 'absolute' poverty line is defined as the minimum income threshold necessary to meet basic needs for survival, such as food, shelter, and clothing, thereby making the statement true.
\end{enumerate}

\section*{Long Form Response}
\begin{enumerate}[label=\arabic*.]
\setcounter{enumi}{10}
\item \textbf{Answer:} In sociology, research purposes typically include exploratory, descriptive, explanatory, and evaluative aims, each serving distinct functions in understanding social phenomena. Exploratory research seeks to investigate new areas where little information exists, while descriptive research focuses on detailing characteristics of a population or phenomenon. Explanatory research aims to identify causal relationships, and evaluative research assesses the effectiveness of programs or interventions. Triangulation, however, challenges these categories as it is not a standalone purpose but rather a methodological approach that integrates multiple data sources or methods to enhance the validity and reliability of research findings. By employing triangulation, researchers can gain a more comprehensive understanding of social issues, but it does not fit neatly into the traditional research purposes because it serves to reinforce or refute findings across existing categories rather than constitute a primary research goal in itself.
\\ \textit{Note:} The answer correctly outlines the four primary research purposes in sociology-exploratory, descriptive, explanatory, and evaluative-by defining each category's specific focus and function, while also implying that triangulation, which involves using multiple methods or data sources, does not fit neatly into these predefined aims.
\item \textbf{Answer:} Monotheism plays a crucial role in shaping the religious organizations of the Church of Norway, Islam, the Church of England, and the Church of Greece, as it establishes a singular focus on one deity that fosters unity and coherence among followers. This shared belief influences their social structures, promoting a sense of community and belonging, as adherents often engage in collective worship and communal activities centered around their faith. Additionally, monotheism impacts cultural practices by guiding moral frameworks, rituals, and traditions that reflect the values and teachings of their respective religions. Consequently, these organizations not only reinforce the authority of a single divine figure but also contribute to the social cohesion and cultural identity of their communities.
\\ \textit{Note:} correct because it identifies monotheism as a unifying force that shapes the social structures and cultural practices of these religious organizations by fostering community engagement and collective worship centered around a singular deity.
\end{enumerate}

\vfill
\noindent\textit{End of Answer Key}
\end{document}